% ============================================================
%  第二卷 第八章 运动电荷的场（§62–§65 + 习题）
%  底本：高教社中文版；OCR 错误已修正，脚注文本待从纸版补录
%
%  OCR 错误修正清单（均已对照原书页面 p-192…p-205 复核）：
%  1. (62.1) 乱码 "$\mathbb{H}$"→"即"；
%  2. §62 正文跨行断句（行8014/8016）"这/些方程的特解" 接合为同段；
%  3. "$\chi(t)=\mathrm{d}\boldsymbol{e}(t)$"→$\chi(t)=\mathrm{d}e(t)$
%     （χ、e 为标量，OCR 误加粗）；
%  4. (62.8) 中 $R=\boldsymbol{r}-\boldsymbol{r}'$、$\boldsymbol{r}=(x,y,z)$
%     等矢量粗体按原书恢复；
%  5. §63 "$t'\mathbb{O}$"→"$t'$①"；习题解 "$\bar\kappa$ 消去"→"从而消去"；
%  6. (63.8) 两项之间的乱码 "$\mathrm{\bf\cdot}$"→"+"（已对照原书 p.182 确认为加号）；
%     E、H 公式按原书分别编号 (63.8)、(63.9)（OCR 的 tag 错位已理顺）；
%  7. §63 "$R^2=R^2$"→"$R^2=\boldsymbol{R}^2$"、
%     "$R\dot R=\boldsymbol{R}\cdot\dot{\boldsymbol{R}}=-\boldsymbol{R}\cdot\boldsymbol{v}$"
%     粗体按原书恢复（p.182、p.188）；
%  8. §64 "$\vec{\boldsymbol{\jmath}}^{*}$ 生场"→"产生场"；习题解中
%     "$(\boldsymbol{\Pi}\varphi)_k$"、"$(\sqcup\varphi)_k$"→"$(\square\varphi)_k$"（达朗贝尔算符）；
%  9. §64 习题 "$\delta(r-vt)$" 等矢量粗体按原书恢复；
% 10. (64.7) 分母 $R(T)$→$R(t)$（原书排印之误，被积变量为 t，已对照原书 p.184；
%     英文原版同式为 R(t)）；
% 11. §65 "已经包含 $\textsf{T}\,1/c$"→"已经包含了 $1/c$"；
%     "（我们已经将 $\dot{\boldsymbol{j}}=\rho\boldsymbol{v}$ 代入）"的
%     $\dot{\boldsymbol{j}}$→$\boldsymbol{j}$（p.187）；
% 12. §65 习题1 "于是我们有： 5" 的孤字"5"为 OCR 噪声，删；
%     "$\boldsymbol{r}_\alpha$"→"$\boldsymbol{r}_a$"；
% 13. §65 习题1 首式 $R$、$\boldsymbol{r}_a$ 及 $\int W\boldsymbol{r}\,\mathrm{d}V$
%     的粗体按原书恢复（p.189）；
% 14. 脚注圈码 ①②③ 按规范保留在正文中，脚注文本未转录（见脚注页码对照.txt
%     行8356/8420/8514 等条目）。
%
%  脚注已转录（2026-10-10）：4 处圈码全部替换为 \footnote{...}
%  （转录自原书扫描页 p-195/200/201/204，与 MinerU OCR page_footnote 块逐字核对）。
%
%  遗留问题：
%  * §62 "这个推导用四维形式为最便利" 照原书（p.178）照排，未改写；
%  * §65 习题1 第二个积分"$\int\nabla(\varphi^2/2)\,\mathrm{d}V$"与第三积分
%    "$\varphi\Delta\varphi\cdot\boldsymbol{r}$"的记号（含点乘）照原书 p.189 照排。
% ============================================================
\chapter{运动电荷的场}

\section{推迟势}

在第五章中我们研究了静止电荷所产生的恒定场，而在第六章中我们研究了不存在电荷的变化场.现在我们来研究存在任意运动电荷的变化场.

我们来推导决定任意运动电荷的场的势的方程.这个推导用四维形式为最便利.重复 §46 末的推导，一个改变是我们将第二对麦克斯韦方程写成（30.2）的形式：
\begin{equation*}
  \frac{\partial F^{ik}}{\partial x^k}=-\frac{4\pi}{c}j^i.
\end{equation*}
右边也同样出现在（46.8）中，在势上施加洛伦兹条件
\begin{equation}
  \frac{\partial A^i}{\partial x^i}=0,\qquad\text{即}\qquad
  \frac{1}{c}\frac{\partial\varphi}{\partial t}+\operatorname{div}\boldsymbol{A}=0,
\end{equation}
以后，我们得到
\begin{equation}
  \frac{\partial^2 A^i}{\partial x_k\partial x^k}=\frac{4\pi}{c}j^i.
\end{equation}

这就是决定任意电磁场势的方程.如用三维形式表示，这个方程可以写成两个方程，一个是 $\boldsymbol{A}$ 的方程，一个是 $\varphi$ 的方程：
\begin{equation}
  \Delta\boldsymbol{A}-\frac{1}{c^2}\frac{\partial^2\boldsymbol{A}}{\partial t^2}
  =-\frac{4\pi}{c}\boldsymbol{j},
\end{equation}
\begin{equation}
  \Delta\varphi-\frac{1}{c^2}\frac{\partial^2\varphi}{\partial t^2}=-4\pi\rho.
\end{equation}

对于恒定场，这些方程就化为熟知的方程（36.4）和（43.4），而对于没有电荷的变化场就化为齐次波动方程.

我们知道，非齐次线性方程（62.3）和（62.4）的解可以用一个和来表示，这个和的第一项是不要右边项的这些方程的通解，第二项是加上右边项的这些方程的特解.为了寻找这个特解，我们将整个空间分为无限小的区域，求其中的一个体积元内的电荷所产生的场.由于这些场方程是线性的，实际的场将是每个体积元所产生的场之和.

在某一体积元内的电荷 $\mathrm{d}e$ 一般来说是时间的函数.假如我们将坐标原点选择在我们所考虑的体积元内，那么，电荷密度是 $\rho=\mathrm{d}e(t)\,\delta(\boldsymbol{R})$，此处的 $\boldsymbol{R}$ 是到原点的距离.因此，我们必须解方程
\begin{equation}
  \Delta\varphi-\frac{1}{c^2}\frac{\partial^2\varphi}{\partial t^2}
  =-4\pi\,\mathrm{d}e(t)\,\delta(\boldsymbol{R}).
\end{equation}

除原点外，到处 $\delta(\boldsymbol{R})=0$，因而我们有方程
\begin{equation}
  \Delta\varphi-\frac{1}{c^2}\frac{\partial^2\varphi}{\partial t^2}=0.
\end{equation}

显然，在我们所考虑的情形下，$\varphi$ 具有中心对称性，就是说，$\varphi$ 仅是 $R$ 的函数.因此，假如我们将拉普拉斯算符写成球坐标形式，（62.6）就化为
\begin{equation*}
  \frac{1}{R^2}\frac{\partial}{\partial R}
  \left(R^2\frac{\partial\varphi}{\partial R}\right)
  -\frac{1}{c^2}\frac{\partial^2\varphi}{\partial t^2}=0.
\end{equation*}

为了解这个方程，我们作 $\varphi=\chi(R,t)/R$ 的替换，那么，对于 $\chi$ 我们得到方程
\begin{equation*}
  \frac{\partial^2\chi}{\partial R^2}
  -\frac{1}{c^2}\frac{\partial^2\chi}{\partial t^2}=0.
\end{equation*}

而这是平面波的方程，它的解有如下的形式（见 §47）：
\begin{equation*}
  \chi=f_1\left(t-\frac{R}{c}\right)+f_2\left(t+\frac{R}{c}\right).
\end{equation*}

因为我们只求方程的特解，因此只选函数 $f_1$ 和 $f_2$ 中的任意一个就够了.通常便利的做法是选取 $f_2=0$（关于这点，可参看后面）.因此，除了原点外，$\varphi$ 到处有如下形式：
\begin{equation}
  \varphi=\frac{\chi(t-R/c)}{R}.
\end{equation}

这个等式中的函数 $\chi$ 暂且是任意的；现在我们如此地选择它，使得在原点也能得到势的正确的值.换句话说，我们必须如此地选择 $\chi$，使方程（62.5）在原点也得到满足.这是很容易办到的，须注意当 $R\to 0$ 时，势将趋向无穷大，因此，它对坐标的导数比它对时间的导数增加得快些.所以，当 $R\to 0$ 时，在方程（62.5）中，同 $\Delta\varphi$ 比较，$\dfrac{1}{c^2}\dfrac{\partial^2\varphi}{\partial t^2}$ 这一项就可以略去了.这时，（62.5）化为熟知的方程（36.9）,后者可导出库仑定律.因此，在原点的附近，（62.7）应该化为库仑定律，由此可以推出 $\chi(t)=\mathrm{d}e(t)$，即
\begin{equation*}
  \varphi=\frac{\mathrm{d}e\left(t-\dfrac{R}{c}\right)}{R}.
\end{equation*}

由此很容易推得方程（62.4）对于任意分布电荷 $\rho(x,y,z,t)$ 的解.为此我们写出 $\mathrm{d}e=\rho\,\mathrm{d}V$（$\mathrm{d}V$ 是体积元），并在整个空间取积分就够了. 在所得的非齐次方程（62.4）的解上，我们还要加上同一方程去掉右边后的解 $\varphi_0$.由此，通解有下面的形式：
\begin{equation}
\begin{aligned}
  &\varphi(\boldsymbol{r},t)=\int\frac{1}{R}\,
  \rho\left(\boldsymbol{r}',t-\frac{R}{c}\right)\mathrm{d}V'+\varphi_0,\\
  &\boldsymbol{R}=\boldsymbol{r}-\boldsymbol{r}',\qquad
  \mathrm{d}V'=\mathrm{d}x'\,\mathrm{d}y'\,\mathrm{d}z',
\end{aligned}
\end{equation}
式中
\begin{equation*}
  \boldsymbol{r}=(x,y,z),\qquad
  \boldsymbol{r}'=(x',y',z');
\end{equation*}
$R$ 是从体积元 $\mathrm{d}V$ 到我们求势的“观察点”的距离.我们将此式简写为
\begin{equation}
  \varphi=\int\frac{\rho_{t-R/c}}{R}\,\mathrm{d}V+\varphi_0,
\end{equation}
式中，下标 $t-R/c$ 表明 $\rho$ 这个量应当取 $t-R/c$ 这一时刻的值，而在 $\mathrm{d}V$ 上的一撇已略去了.

对于矢势我们类似地得到：
\begin{equation}
  \boldsymbol{A}=\frac{1}{c}\int\frac{\boldsymbol{j}_{t-R/c}}{R}\,\mathrm{d}V
  +\boldsymbol{A}_0,
\end{equation}
式中，$\boldsymbol{A}_0$ 是方程（62.3）去掉右边项的解.

（62.9）和（62.10）式（不带 $\varphi_0$ 和 $\boldsymbol{A}_0$）称为推迟势.

假如电荷是静止的（即 $\rho$ 与时间无关），公式（62.9）就化为熟知的静电场的公式（36.8）；而对于电荷平稳运动的情形（62.10）经过求平均以后则变为恒定磁场的公式（43.5）.

通过决定（62.9）式中的 $\varphi_0$ 和（62.10）式中的 $\boldsymbol{A}_0$，要使得问题的条件能够得到满足.为此，显然只需给出初始条件，即给出场在开始时的值就够了.但是通常我们并不必须讨论这些初始条件.作为代替，我们给出在与电荷系统相距很远之处、在所有时刻的条件.也就是说，辐射是从系统外面入射到系统上的.与此相应，辐射与系统相互作用的结果所建立的场与外场的差别仅仅是系统所发出的辐射.这个系统所发出的辐射在远处有波的形式，由系统向外沿着 $R$ 增加的方向传播.但是这个条件正好为推迟势所满足.因此，这些解代表系统所产生的场，而 $\varphi_0$ 和 $\boldsymbol{A}_0$ 必须等于作用于系统上的外场.

\section{李纳-维谢尔势}

一个点电荷沿轨道 $\boldsymbol{r}=\boldsymbol{r}_0(t)$ 进行给定的运动，让我们来求它所产生的场的势.

按照推迟势的公式，时刻 $t$ 在观察点 $P(x,y,z)$ 的场，由电荷在较早时刻 $t'$ 的运动状态决定，光信号从电荷所在点 $\boldsymbol{r}_0(t')$ 传播到观察点 $P$ 的时间正好与差 $t-t'$ 一致.令 $\boldsymbol{R}(t)=\boldsymbol{r}-\boldsymbol{r}_0(t)$ 为电荷 $e$ 到点 $P$ 的径矢；像 $\boldsymbol{r}_0(t)$ 一样，它是时间的给定函数.于是时刻 $t'$ 由如下方程决定
\begin{equation}
  t'+\frac{R(t')}{c}=t.
\end{equation}

对于每个 $t$ 值，这个方程只有一个根 $t'$\footnote{这一点是显然的，但可以直接验证.为此，我们选择观察点 $P$ 和观察时刻 $t$ 作为四维坐标系的原点 $O$，并以 $O$ 为其顶点构造光锥（§2）.这个锥的下半部分包围着“绝对”过去（对事件 $O$ 而言）的区域，是这样一些世界点构成的几何图形，光信号可从这些点达到点 $O$.这个超曲面同电荷的世界线的交点显然正是（63.1）的根.但是因为粒子的速度总小于光速，它的运动的世界线对于时轴的斜度总处处小于光锥的斜度.由此可知，粒子的世界线仅仅在一点能与光锥的下半部分相交.}.

在 $t'$ 时刻粒子为静止的参考系中，观察点在 $t$ 时刻的势正好是库仑势，即
\begin{equation}
  \varphi=\frac{e}{R(t')},\qquad \boldsymbol{A}=0.
\end{equation}

在任意的参考系中，势的表达式可以直接从求一个四维矢量得到，这个四维矢量当 $v=0$ 时与方才所得的 $\varphi$ 及 $\boldsymbol{A}$ 的表达式相合.应当注意，按照（63.1），（63.2）中的 $\varphi$ 也可以写成下面的形式：
\begin{equation*}
  \varphi=\frac{e}{c(t-t')},
\end{equation*}
我们得到所要求的四维矢量是：
\begin{equation}
  A^i=e\,\frac{u^i}{R_ku^k},
\end{equation}
其中 $u^k$ 是电荷的四维速度，而 $R^k=[c(t-t'),\boldsymbol{r}-\boldsymbol{r}']$，式中 $x',y',z',t'$ 通过（63.1）彼此相关；写成四维形式是
\begin{equation}
  R_kR^k=0.
\end{equation}

现在再变为三维表示法，我们得到一个任意运动的点电荷所产生的场的势如下：
\begin{equation}
  \varphi=\frac{e}{\left(R-\dfrac{\boldsymbol{v}\cdot\boldsymbol{R}}{c}\right)},\qquad
  \boldsymbol{A}=\frac{e\,\boldsymbol{v}}{c\left(R-\dfrac{\boldsymbol{v}\cdot\boldsymbol{R}}{c}\right)},
\end{equation}
其中 $\boldsymbol{R}$ 是从电荷所在点到观察点 $P$ 的径矢，而在方程右边的所有的量都必须取在（63.1）所决定时刻 $t'$ 的值.场的势写成（63.5）的形式称为李纳-维谢尔势.

为了从公式
\begin{equation*}
  \boldsymbol{E}=-\frac{1}{c}\frac{\partial\boldsymbol{A}}{\partial t}
  -\operatorname{grad}\varphi,\qquad
  \boldsymbol{H}=\operatorname{rot}\boldsymbol{A}
\end{equation*}
来计算电场和磁场的强度，我们必须将 $\varphi$ 及 $\boldsymbol{A}$ 对点的坐标 $x,y,z$ 和观察时刻 $t$ 微分.但是公式（63.5)是用 $t'$ 的函数来表示势，而只有通过关系式（63.1）才能表示势为 $x,y,z,t$ 的隐函数.因此，为了计算所需要的导数，必须先计算对 $t'$ 的导数.将关系式 $R(t')=c(t-t')$ 对 $t$ 微分，得
\begin{equation*}
  \frac{\partial R}{\partial t}
  =\frac{\partial R}{\partial t'}\frac{\partial t'}{\partial t}
  =-\frac{\boldsymbol{R}\cdot\boldsymbol{v}}{R}\frac{\partial t'}{\partial t}
  =c\left(1-\frac{\partial t'}{\partial t}\right).
\end{equation*}
（$\partial R/\partial t'$ 的值由微分恒等式 $R^2=\boldsymbol{R}^2$，并将
$\partial\boldsymbol{R}(t')/\partial t'=-\boldsymbol{v}(t')$ 代入来求得.前面有负号是因为 $\boldsymbol{R}$ 是从电荷 $e$ 到 $P$ 点的径矢，而不是从 $P$ 点到电荷 $e$ 的径矢.）因此，
\begin{equation}
  \frac{\partial t'}{\partial t}
  =\frac{1}{1-\dfrac{\boldsymbol{v}\cdot\boldsymbol{R}}{Rc}}.
\end{equation}

同理，将同一个关系式对坐标微分，我们得到
\begin{equation*}
  \operatorname{grad}t'=-\frac{1}{c}\operatorname{grad}R(t')
  =-\frac{1}{c}\left(\frac{\partial R}{\partial t'}\operatorname{grad}t'
  +\frac{\boldsymbol{R}}{R}\right),
\end{equation*}
所以
\begin{equation}
  \operatorname{grad}t'
  =-\frac{\boldsymbol{R}}{c\left(R-\dfrac{\boldsymbol{R}\cdot\boldsymbol{v}}{c}\right)}.
\end{equation}

利用这些公式，求场 $\boldsymbol{E}$ 及 $\boldsymbol{H}$ 的运算过程就不难了.略去运算的中间过程，我们得到结果：
\begin{equation}
\begin{aligned}
  \boldsymbol{E}=e\,\frac{1-\dfrac{v^2}{c^2}}
  {\left(R-\dfrac{\boldsymbol{R}\cdot\boldsymbol{v}}{c}\right)^{3}}
  &\left(\boldsymbol{R}-\frac{\boldsymbol{v}}{c}R\right)\\
  &+\frac{e}{c^2\left(R-\dfrac{\boldsymbol{R}\cdot\boldsymbol{v}}{c}\right)^{3}}
  \,\boldsymbol{R}\times\left\{\left(\boldsymbol{R}
  -\frac{\boldsymbol{v}}{c}R\right)\times\dot{\boldsymbol{v}}\right\},
\end{aligned}
\end{equation}
\begin{equation}
  \boldsymbol{H}=\frac{1}{R}\,\boldsymbol{R}\times\boldsymbol{E}.
\end{equation}

式中 $\dot{\boldsymbol{v}}=\partial\boldsymbol{v}/\partial t'$；等式右边所有的量都是取在 $t'$ 时刻的值.值得注意，磁场无论在何处都与电场垂直.

电场（63.8）由两个不同类型的部分组成.第一项只依赖于粒子的速度（而不依赖于其加速度），且在大距离处像 $1/R^2$ 那样变化.第二项依赖于加速度，且对于大的 $R$ 像 $1/R$ 那样变化.以后（§66）我们将看到，这个第二项与粒子辐射的电磁波有关.

至于第一项，因为它与加速度无关，必定对应着匀速运动电荷产生的场.实际上，由于速度恒定，差
\begin{equation*}
  \boldsymbol{R}_{t'}-\frac{\boldsymbol{v}}{c}R_{t'}
  =\boldsymbol{R}_{t'}-\boldsymbol{v}(t-t')
\end{equation*}
就是在准确的观察时刻从电荷到观察点的距离 $\boldsymbol{R}_t$.也容易直接证明，
\begin{equation*}
  R_{t'}-\frac{1}{c}\boldsymbol{R}_{t'}\cdot\boldsymbol{v}
  =\sqrt{R_t^2-\frac{1}{c^2}(\boldsymbol{v}\times\boldsymbol{R}_t)^2}
  =R_t\sqrt{1-\frac{v^2}{c^2}\sin^2\theta_t},
\end{equation*}
式中，$\theta_t$ 是 $\boldsymbol{R}_t$ 和 $\boldsymbol{v}$ 之间的夹角.因而，（63.8）的第一项与表达式（38.8）相同.

\begin{xiti}

\noindent{\heiti 习\ 题}\quad 通过积分（62.9）—（62.10）推导李纳-维谢尔势.

解：我们将（62.8）写为形式
\begin{equation*}
  \varphi(\boldsymbol{r},t)=\iint
  \frac{\rho(\boldsymbol{r}',\tau)}{|\boldsymbol{r}-\boldsymbol{r}'|}
  \,\delta\left(\tau-t+\frac{1}{c}|\boldsymbol{r}-\boldsymbol{r}'|\right)
  \mathrm{d}\tau\,\mathrm{d}V'
\end{equation*}
（对 $\boldsymbol{A}(\boldsymbol{r},t)$ 作同样处理），引入附加的 $\delta$ 函数从而消去函数 $\rho$ 中的隐宗量.对于沿轨道 $\boldsymbol{r}=\boldsymbol{r}_0(t)$ 运动的点电荷我们有：
\begin{equation*}
  \rho(\boldsymbol{r}',\tau)=e\,\delta[\boldsymbol{r}'-\boldsymbol{r}_0(\tau)].
\end{equation*}

将这个表达式代入并对 $\mathrm{d}V'$ 积分，我们得到：
\begin{equation*}
  \varphi(\boldsymbol{r},t)=e\int
  \frac{\mathrm{d}\tau}{|\boldsymbol{r}-\boldsymbol{r}_0(\tau)|}\,
  \delta\left[\tau-t+\frac{1}{c}|\boldsymbol{r}-\boldsymbol{r}_0(\tau)|\right],
\end{equation*}
对 $\tau$ 积分是用如下公式进行的
\begin{equation*}
  \delta[F(\tau)]=\frac{\delta(\tau-t')}{|F'(t')|}
\end{equation*}
（式中 $t'$ 是 $F(t')=0$ 的根），于是得到公式（63.5）.

\end{xiti}

\section{推迟势的谱分解}

运动电荷所产生的场可以展开为单色波.场的不同的单色分量（傅里叶分量）的势具有 $\varphi_\omega\mathrm{e}^{-\mathrm{i}\omega t}$，$\boldsymbol{A}_\omega\mathrm{e}^{-\mathrm{i}\omega t}$ 的形式.产生场的电荷系统的电荷密度和电流密度也都可以展开为傅里叶级数或傅里叶积分.很清楚，$\rho$ 和 $\boldsymbol{j}$ 的傅里叶分量产生场的相应的单色分量.

为了用电荷密度和电流密度的傅里叶分量表示场的傅里叶分量，我们分别用 $\varphi_\omega\mathrm{e}^{-\mathrm{i}\omega t}$ 和 $\rho_\omega\mathrm{e}^{-\mathrm{i}\omega t}$ 代替（62.9）式中的 $\varphi$ 和 $\rho$，这样，就得到
\begin{equation*}
  \varphi_\omega\mathrm{e}^{-\mathrm{i}\omega t}
  =\int\rho_\omega\,\frac{\mathrm{e}^{-\mathrm{i}\omega(t-R/c)}}{R}\,\mathrm{d}V.
\end{equation*}
将 $\mathrm{e}^{-\mathrm{i}\omega t}$ 消去，并引用波矢的绝对值 $k=\omega/c$，我们得到：
\begin{equation}
  \varphi_\omega=\int\rho_\omega\,\frac{\mathrm{e}^{\mathrm{i}kR}}{R}\,\mathrm{d}V.
\end{equation}

同样，对于 $\boldsymbol{A}_\omega$，我们得到
\begin{equation}
  \boldsymbol{A}_\omega=\int\boldsymbol{j}_\omega\,
  \frac{\mathrm{e}^{\mathrm{i}kR}}{cR}\,\mathrm{d}V.
\end{equation}

我们指出，公式（64.1）是更广泛形式的泊松方程
\begin{equation}
  \Delta\varphi_\omega+k^2\varphi_\omega=-4\pi\rho_\omega
\end{equation}
的解（该方程是从（62.4）对于 $\rho$ 和 $\varphi$ 与时间的关系为 $\mathrm{e}^{-\mathrm{i}\omega t}$ 的情形得来的）.

如果我们所研究的是傅里叶积分的展开式，那么，电荷密度的傅里叶分量是
\begin{equation*}
  \rho_\omega=\int_{-\infty}^{+\infty}\rho\,\mathrm{e}^{\mathrm{i}\omega t}\,\mathrm{d}t.
\end{equation*}

将上式代入（64.1），我们得到
\begin{equation}
  \varphi_\omega=\int_{-\infty}^{+\infty}\int
  \frac{\rho}{R}\,\mathrm{e}^{\mathrm{i}(\omega t+kR)}\,\mathrm{d}V\,\mathrm{d}t.
\end{equation}

我们还必须从电荷密度的连续分布过渡到正在考虑其运动的点电荷.因此，若只有一个点电荷，令
\begin{equation*}
  \rho=e\,\delta[\boldsymbol{r}-\boldsymbol{r}_0(t)],
\end{equation*}
式中 $\boldsymbol{r}_0(t)$ 是电荷的径矢，它是时间的已知函数.将此式代入（64.4）再对空间积分（这样积分归结为用 $\boldsymbol{r}_0(t)$ 代替 $\boldsymbol{r}$），我们得到
\begin{equation}
  \varphi_\omega=e\int_{-\infty}^{\infty}\frac{1}{R(t)}\,
  \mathrm{e}^{\mathrm{i}\omega[t+R(t)/c]}\,\mathrm{d}t,
\end{equation}
这里 $R(t)$ 是从电荷到观察点的距离.同样，对于矢势，我们得到
\begin{equation}
  \boldsymbol{A}_\omega=\frac{e}{c}\int_{-\infty}^{\infty}
  \frac{\boldsymbol{v}(t)}{R(t)}\,
  \mathrm{e}^{\mathrm{i}\omega[t+R(t)/c]}\,\mathrm{d}t,
\end{equation}
式中 $\boldsymbol{v}=\dot{\boldsymbol{r}}_0(t)$ 是电荷的速度.

对于电荷密度和电流密度的谱分解含一系列离散频率的情形，也可以写出类似（64.5），（64.6）的公式.因此，对于点电荷的周期运动（周期 $T=2\pi/\omega_0$），场的谱分解只含形如 $n\omega_0$ 的频率，相应的矢势分量是
\begin{equation}
  \boldsymbol{A}_n=\frac{e}{cT}\int_0^{T}
  \frac{\boldsymbol{v}(t)}{R(t)}\,
  \mathrm{e}^{\mathrm{i}n\omega_0[t+R(t)/c]}\,\mathrm{d}t
\end{equation}
（对于 $\varphi_n$ 有类似公式）. 在（64.6)，(64.7)中傅里叶分量都是按 §49 定义的.

\begin{xiti}

\noindent{\heiti 习\ 题}\quad 把匀速直线运动的电荷所产生的场展开为平面波.

解：我们按照 §51 中的方式进行.将电荷密度写为 $\rho=e\,\delta(\boldsymbol{r}-\boldsymbol{v}t)$ 的形式，此处的 $\boldsymbol{v}$ 是电荷的速度.取方程
$\square\varphi=-4\pi e\,\delta(\boldsymbol{r}-\boldsymbol{v}t)$ 的傅里叶分量，我们得到
\begin{equation*}
  (\square\varphi)_{\boldsymbol{k}}
  =-4\pi e\cdot\mathrm{e}^{-\mathrm{i}(\boldsymbol{v}\cdot\boldsymbol{k})t}.
\end{equation*}

另一方面，由于
\begin{equation*}
  \varphi=\int\mathrm{e}^{\mathrm{i}\boldsymbol{k}\cdot\boldsymbol{r}}\,
  \varphi_{\boldsymbol{k}}\,\frac{\mathrm{d}^3k}{(2\pi)^3}
\end{equation*}
所以我们有
\begin{equation*}
  (\square\varphi)_{\boldsymbol{k}}
  =-k^2\varphi_{\boldsymbol{k}}
  -\frac{1}{c^2}\frac{\partial^2\varphi_{\boldsymbol{k}}}{\partial t^2}.
\end{equation*}

因此，
\begin{equation*}
  \frac{1}{c^2}\frac{\partial^2\varphi_{\boldsymbol{k}}}{\partial t^2}
  +k^2\varphi_{\boldsymbol{k}}
  =4\pi e\cdot\mathrm{e}^{-\mathrm{i}(\boldsymbol{k}\cdot\boldsymbol{v})t},
\end{equation*}
从而，最后得到
\begin{equation*}
  \varphi_{\boldsymbol{k}}=4\pi e\,
  \frac{\mathrm{e}^{-\mathrm{i}(\boldsymbol{k}\cdot\boldsymbol{v})t}}
       {k^2-\left(\dfrac{\boldsymbol{k}\cdot\boldsymbol{v}}{c}\right)^2}.
\end{equation*}

从此可以推断，以 $\boldsymbol{k}$ 为波矢量的波的频率是 $\omega=\boldsymbol{k}\cdot\boldsymbol{v}$.

同样，我们得到矢势：
\begin{equation*}
  \boldsymbol{A}_{\boldsymbol{k}}=\frac{4\pi e}{c}\,
  \frac{\boldsymbol{v}\,\mathrm{e}^{-\mathrm{i}(\boldsymbol{k}\cdot\boldsymbol{v})t}}
       {k^2-\left(\dfrac{\boldsymbol{k}\cdot\boldsymbol{v}}{c}\right)^2}.
\end{equation*}

最后，我们得到场的表达式如下：
\begin{equation*}
  \boldsymbol{E}_{\boldsymbol{k}}
  =-\mathrm{i}\boldsymbol{k}\,\varphi_{\boldsymbol{k}}
  +\mathrm{i}\,\frac{\boldsymbol{k}\cdot\boldsymbol{v}}{c}\,
  \boldsymbol{A}_{\boldsymbol{k}}
  =\mathrm{i}\,4\pi e\,
  \frac{-\boldsymbol{k}+\dfrac{(\boldsymbol{k}\cdot\boldsymbol{v})}{c^2}\boldsymbol{v}}
       {k^2-\left(\dfrac{\boldsymbol{k}\cdot\boldsymbol{v}}{c}\right)^2}
  \,\mathrm{e}^{-\mathrm{i}(\boldsymbol{k}\cdot\boldsymbol{v})t},
\end{equation*}
\begin{equation*}
  \boldsymbol{H}_{\boldsymbol{k}}
  =\mathrm{i}\boldsymbol{k}\times\boldsymbol{A}_{\boldsymbol{k}}
  =\mathrm{i}\,\frac{4\pi e}{c}\,
  \frac{\boldsymbol{k}\times\boldsymbol{v}}
       {k^2-\left(\dfrac{\boldsymbol{k}\cdot\boldsymbol{v}}{c}\right)^2}
  \,\mathrm{e}^{-\mathrm{i}(\boldsymbol{k}\cdot\boldsymbol{v})t}.
\end{equation*}

\end{xiti}

\section{精确到二阶的拉格朗日量}

在普通经典力学中，我们利用仅仅与粒子（在同一时刻）的坐标和速度有关的拉格朗日量来描写彼此相互作用的粒子体系.这种作法的可行性，归根到底是来自于，在经典力学中假定粒子间的相互作用的传播速度是无穷大.

我们已经知道，因为相互作用的传播速度是有限的，场必须被看做是一个有本身“自由度”的独立体系.由此可得出结论：假如有一个相互作用的粒子（电荷）体系，那么，为了描述它，我们就必须认为这个体系由这些粒子和场组成.因此，当考虑到相互作用传播速度的有限性时，一般说来，不可能用仅仅和粒子的坐标和速度有关，而并不包含与场的内部“自由度”有关量的拉格朗日量来严格描写相互作用的粒子体系.

然而，如果所有粒子的速度 $v$ 都比光速小很多，那么，这个电荷体系就可以用某个近似的拉格朗日量来描写.这时，就可能引入一个拉格朗日量，这个拉格朗日量不仅在忽略所有 $v/c$ 的幂的情况下描写这个体系（经典的拉格朗日量），而且还准确到 $v^2/c^2$ 的数量级.最后这个论断与如下事实有关，即运动电荷辐射电磁波（因而出现“自”场)，只是在 $v/c$ 的三级近似下发生（参见下面的 §67）\footnote{对于由荷质比相同的粒子组成的系统，辐射在 $v/c$ 的第五级近似才发生；在这种情形下，拉格朗日量含有直到 $v/c$ 的四级项.（参见 \textit{Barker B. M., O'Connel R. F.}//Canad. J. Phys. 1980. V. 58. P. 1659）.}.

我们预先指出，在零级近似中，即当我们完全略去势的推迟时，电荷体系的拉格朗日量有以下的形式：
\begin{equation}
  L^{(0)}=\sum_a\frac{m_av_a^2}{2}
  -\sum_{a>b}\frac{e_ae_b}{R_{ab}}
\end{equation}
（求和应该对于组成体系的所有电荷进行）.第二项是相互作用的势能，正如静止电荷的情况一样.

为了得到拉格朗日量的高一级的近似式，我们按下述方式进行.电荷 $e_a$ 在外场中的拉格朗日量是
\begin{equation}
  L_a=-m_ac^2\sqrt{1-\frac{v_a^2}{c^2}}
  -e_a\varphi+\frac{e_a}{c}\,\boldsymbol{A}\cdot\boldsymbol{v}_a.
\end{equation}

选定电荷体系中的任意一个电荷，我们求出所有其余电荷在第一电荷所在点产生的场的势，并且用产生这个场的电荷的坐标和速度来表示它们（这一点只能近似作到——对于 $\varphi$，精确到 $v^2/c^2$；对于 $\boldsymbol{A}$，精确到 $v/c$）.将这样求得的势的表达式代入（65.2），我们就得到电荷体系中的一个电荷的拉格朗日量（当其余电荷的运动已知时）.从此，不难求出整个体系的拉格朗日量 $L$.

我们从推迟势的表达式：
\begin{equation*}
  \varphi=\int\frac{\rho_{t-R/c}}{R}\,\mathrm{d}V,\qquad
  \boldsymbol{A}=\frac{1}{c}\int\frac{\boldsymbol{j}_{t-R/c}}{R}\,\mathrm{d}V
\end{equation*}
出发，如果所有电荷的速度都比光速小很多，那么，电荷的分布在时间 $R/c$ 内不致于有显著的改变.因此，我们可以将 $\rho_{t-R/c}$ 和 $\boldsymbol{j}_{t-R/c}$ 展开为 $R/c$ 的幂级数.因此，对于标势，我们得到精确到二阶项的表达式
\begin{equation*}
  \varphi=\int\frac{\rho\,\mathrm{d}V}{R}
  -\frac{1}{c}\frac{\partial}{\partial t}\int\rho\,\mathrm{d}V
  +\frac{1}{2c^2}\frac{\partial^2}{\partial t^2}\int R\rho\,\mathrm{d}V
\end{equation*}
（没有下标的 $\rho$ 是指它在 $t$ 时刻的值，$\dfrac{\partial}{\partial t}$ 和
$\dfrac{\partial^2}{\partial t^2}$ 显然可以从积分号内提出）.但是 $\int\rho\,\mathrm{d}V$ 是总电荷，它是一个与时间无关的常量.因此上式中的第二项为零，从而
\begin{equation}
  \varphi=\int\frac{\rho\,\mathrm{d}V}{R}
  +\frac{1}{2c^2}\frac{\partial^2}{\partial t^2}\int R\rho\,\mathrm{d}V.
\end{equation}

对于 $\boldsymbol{A}$，我们也可以用同样方式进行.但是用电流密度表示矢势的式子已经包含了 $1/c$，而当我们把它代入拉格朗日量时，还要乘上 $1/c$. 既然我们正在求仅仅精确到二阶的拉格朗日量，我们在 $\boldsymbol{A}$ 的展开式中只取其第一项就足够了，即
\begin{equation}
  \boldsymbol{A}=\frac{1}{c}\int\frac{\rho\boldsymbol{v}}{R}\,\mathrm{d}V
\end{equation}
（我们已经将 $\boldsymbol{j}=\rho\boldsymbol{v}$ 代入）.

假设只有一个单独的点电荷 $e$.我们从（65.3）和（65.4）两式得到
\begin{equation}
  \varphi=\frac{e}{R}+\frac{e}{2c^2}\frac{\partial^2 R}{\partial t^2},\qquad
  \boldsymbol{A}=\frac{e\,\boldsymbol{v}}{cR},
\end{equation}
式中 $R$ 是到电荷的距离.

通过变换
\begin{equation*}
  \varphi'=\varphi-\frac{1}{c}\frac{\partial f}{\partial t},\qquad
  \boldsymbol{A}'=\boldsymbol{A}+\operatorname{grad}f,
\end{equation*}
我们选择另外两个势 $\varphi'$ 和 $\boldsymbol{A}'$ 来代替 $\varphi$ 和 $\boldsymbol{A}$（§18）.并且我们选择 $f$ 如下：
\begin{equation*}
  f=\frac{e}{2c}\,\frac{\partial R}{\partial t}.
\end{equation*}

这时，我们得到\footnote{这些势不再满足洛伦兹条件（62.1），也不满足（62.3）—（62.4）.}
\begin{equation*}
  \varphi'=\frac{e}{R},\qquad
  \boldsymbol{A}'=\frac{e\,\boldsymbol{v}}{cR}
  +\frac{e}{2c}\,\nabla\frac{\partial R}{\partial t}.
\end{equation*}

为了计算 $\boldsymbol{A}'$，首先我们注意到
$\nabla\dfrac{\partial R}{\partial t}=\dfrac{\partial}{\partial t}\nabla R$.在此处的符号 $\nabla$ 是指对我们正在求 $\boldsymbol{A}'$ 值的观察点坐标微分.因此 $\nabla R$ 是一个单位矢量 $\boldsymbol{n}$，从电荷 $e$ 指向观察点，所以
\begin{equation*}
  \boldsymbol{A}'=\frac{e\,\boldsymbol{v}}{cR}+\frac{e}{2c}\,\dot{\boldsymbol{n}}.
\end{equation*}

为了计算 $\dot{\boldsymbol{n}}$，我们写出
\begin{equation*}
  \dot{\boldsymbol{n}}=\frac{\partial}{\partial t}
  \left(\frac{\boldsymbol{R}}{R}\right)
  =\frac{\dot{\boldsymbol{R}}}{R}-\frac{R\dot{R}}{R^2}.
\end{equation*}

但是对于一指定的观察点，导数 $\dot{R}$ 是电荷速度 $\boldsymbol{v}$，且只要微分等式 $R^2=\boldsymbol{R}^2$，亦即，写出
\begin{equation*}
  R\dot{R}=\boldsymbol{R}\cdot\dot{\boldsymbol{R}}=-\boldsymbol{R}\cdot\boldsymbol{v}
\end{equation*}
就很容易确定导数 $\dot{R}$.因此，
\begin{equation*}
  \dot{\boldsymbol{n}}=\frac{-\boldsymbol{v}+\boldsymbol{n}(\boldsymbol{n}\cdot\boldsymbol{v})}{R}.
\end{equation*}

将此式代入 $\boldsymbol{A}'$ 的表达式，我们最后得到
\begin{equation}
  \varphi'=\frac{e}{R},\qquad
  \boldsymbol{A}'=\frac{e[\boldsymbol{v}+(\boldsymbol{v}\cdot\boldsymbol{n})\boldsymbol{n}]}{2cR}.
\end{equation}

如果有许多电荷，则显然必须对所有的电荷求和.

若将这些势的表达式代入（65.2），我们就得到一个电荷 $e_a$ 的拉格朗日量 $L_a$（当其余所有电荷的运动为已知时）.这时，我们也必须将（65.2)的第一项展开为 $v_a/c$ 的幂级数，并保留到二阶项为止.这样，我们就得到如下 $L_a$ 的表达式：
\begin{equation*}
\begin{aligned}
  L_a={}&\frac{m_av_a^2}{2}+\frac{1}{8}\,\frac{m_av_a^4}{c^2}
  -e_a\sum_b'\frac{e_b}{R_{ab}}\\
  &+\frac{e_a}{2c^2}\sum_b'\frac{e_b}{R_{ab}}
  \bigl[\boldsymbol{v}_a\cdot\boldsymbol{v}_b
  +(\boldsymbol{v}_a\cdot\boldsymbol{n}_{ab})
  (\boldsymbol{v}_b\cdot\boldsymbol{n}_{ab})\bigr]
\end{aligned}
\end{equation*}
（求和应对除 $e_a$ 以外的其余所有电荷进行；$\boldsymbol{n}_{ab}$ 是从 $e_b$ 到 $e_a$ 的单位矢量）.

由此就不难求出整个体系的拉格朗日量.不难理解，这个函数不是所有电荷的 $L_a$ 之和，而有如下的形式：
\begin{equation}
\begin{aligned}
  L={}&\sum_a\frac{m_av_a^2}{2}+\sum_a\frac{m_av_a^4}{8c^2}
  -\sum_{a>b}\frac{e_ae_b}{R_{ab}}\\
  &+\sum_{a>b}\frac{e_ae_b}{2c^2R_{ab}}
  \bigl[\boldsymbol{v}_a\cdot\boldsymbol{v}_b
  +(\boldsymbol{v}_a\cdot\boldsymbol{n}_{ab})
  (\boldsymbol{v}_b\cdot\boldsymbol{n}_{ab})\bigr].
\end{aligned}
\end{equation}

实际上，对于电荷中的每一个，在其余所有电荷的运动为已知的情况下，函数 $L$ 化为上面求得的 $L_a$.(65.7) 式决定电荷体系精确到二阶项的拉格朗日量（由 C. G. Darwin，1922 首先得到）.

最后，我们还要求电荷体系的哈密顿量，并要有同样的近似程度，这可应用从 $L$ 求 $\mathcal{H}$ 的一般法则进行；然而用下面的方法来求则更为简单.在（65.7）式中的第二和第四两项是对 $L^{(0)}$ (65.1) 的小修正.另一方面，从力学我们知道当 $L$ 和 $\mathcal{H}$ 有小的变化时，加到它们上面的小量的数值相等，正负号相反（在此，$L$ 的变化是发生在坐标和速度恒定的情况下，而 $\mathcal{H}$ 的变化，是发生在坐标和动量恒定的情况下；参见本教程第一卷 §40）.

因此我们从
\begin{equation*}
  \mathcal{H}^{(0)}=\sum_a\frac{p_a^2}{2m_a}
  +\sum_{a>b}\frac{e_ae_b}{R_{ab}}
\end{equation*}
中减去（65.7）式的第二和第四两项，用一级近似 $\boldsymbol{v}_a=\boldsymbol{p}_a/m_a$ 代替其中的速度，就可立刻写出 $\mathcal{H}$，所以
\begin{equation}
\begin{aligned}
  \mathcal{H}={}&\sum_a\frac{p_a^2}{2m_a}
  -\sum_a\frac{p_a^4}{8c^2m_a^3}
  +\sum_{a>b}\frac{e_ae_b}{R_{ab}}\\
  &-\sum_{a>b}\frac{e_ae_b}{2c^2m_am_bR_{ab}}
  \bigl[\boldsymbol{p}_a\cdot\boldsymbol{p}_b
  +(\boldsymbol{p}_a\cdot\boldsymbol{n}_{ab})
  (\boldsymbol{p}_b\cdot\boldsymbol{n}_{ab})\bigr].
\end{aligned}
\end{equation}

\begin{xiti}

\noindent{\heiti 习题1}\quad 确定相互作用的粒子体系的惯性中心（精确到二阶项）.

解：解决这个问题最简单的办法是用公式
\begin{equation*}
  \boldsymbol{R}=\frac{\displaystyle\sum_a\mathcal{E}_a\boldsymbol{r}_a
  +\int W\boldsymbol{r}\,\mathrm{d}V}
  {\displaystyle\sum_a\mathcal{E}_a+\int W\,\mathrm{d}V}
\end{equation*}
（参见（14.6)），式中 $\mathcal{E}_a$ 是粒子的动能（包括它的静能），$W$ 是粒子产生的场的能量密度.因为 $\mathcal{E}_a$ 包含着很大的量 $m_ac^2$，在得到下一级近似时，只考虑 $\mathcal{E}_a$ 和 $W$ 中那些不含 $c$ 的项就够了.也就是说，我们只需要考虑粒子的非相对论动能和静电场的能量.于是我们有：
\begin{equation*}
\begin{aligned}
  \int W\boldsymbol{r}\,\mathrm{d}V
  &=\frac{1}{8\pi}\int E^2\boldsymbol{r}\,\mathrm{d}V
  =\frac{1}{8\pi}\int(\nabla\varphi)^2\boldsymbol{r}\,\mathrm{d}V=\\
  &=\frac{1}{8\pi}\int\left(\mathrm{d}\boldsymbol{f}\cdot\nabla\frac{\varphi^2}{2}\right)
  \boldsymbol{r}
  -\frac{1}{8\pi}\int\nabla\frac{\varphi^2}{2}\,\mathrm{d}V
  -\frac{1}{8\pi}\int\varphi\Delta\varphi\cdot\boldsymbol{r}\,\mathrm{d}V;
\end{aligned}
\end{equation*}
无穷远表面的积分为零；第二个积分也换成面积分变为零，而在第三个积分中代入 $\Delta\varphi=-4\pi\rho$ 后得：
\begin{equation*}
  \int W\boldsymbol{r}\,\mathrm{d}V
  =\frac{1}{2}\int\rho\varphi\,\boldsymbol{r}\,\mathrm{d}V
  =\frac{1}{2}\sum_a e_a\varphi_a\boldsymbol{r}_a,
\end{equation*}
式中 $\varphi_a$ 是除 $e_a$ 外所有的电荷在点 $\boldsymbol{r}_a$ 产生的势\footnote{去掉粒子的自场相应于 §37 第 1 个脚注提到的质量“重正化”.}.

最后，我们得到：
\begin{equation*}
  \boldsymbol{R}=\frac{1}{\mathcal{E}}\sum_a\boldsymbol{r}_a
  \left(m_ac^2+\frac{p_a^2}{2m_a}
  +\frac{e_a}{2}\sum_b'\frac{e_b}{R_{ab}}\right)
\end{equation*}
（求和遍历除 $b=a$ 以外的所有指标 $b$），式中
\begin{equation*}
  \mathcal{E}=\sum_a\left(m_ac^2+\frac{p_a^2}{2m_a}
  +\sum_{a>b}\frac{e_ae_b}{R_{ab}}\right)
\end{equation*}
是系统的总能量.因此在这个近似下，惯性中心的坐标实际上可以用只针对粒子的量来表示.

\noindent{\heiti 习题2}\quad 试求由两个粒子构成体系的哈密顿量（精确到二阶，略去体系整体的运动）.

解：我们选择一个参考系，其中两个粒子的总动量为零.将动量写为作用量的导数，则我们有：
\begin{equation*}
  \boldsymbol{p}_1+\boldsymbol{p}_2
  =\frac{\partial S}{\partial\boldsymbol{r}_1}
  +\frac{\partial S}{\partial\boldsymbol{r}_2}=0.
\end{equation*}

由此可见，在我们所选择的参考系中，作用量是两个粒子径矢之差
$\boldsymbol{r}=\boldsymbol{r}_2-\boldsymbol{r}_1$ 的函数，因此，我们有
$\boldsymbol{p}_2=-\boldsymbol{p}_1=\boldsymbol{p}$，这里的
$\boldsymbol{p}=\partial S/\partial\boldsymbol{r}$ 是两个粒子的相对运动的动量.哈密顿量等于
\begin{equation*}
  \mathcal{H}=\frac{p^2}{2}\left(\frac{1}{m_1}+\frac{1}{m_2}\right)
  +\frac{e_1e_2}{r}
  -\frac{p^4}{8c^2}\left(\frac{1}{m_1^3}+\frac{1}{m_2^3}\right)
  +\frac{e_1e_2}{2m_1m_2c^2r}\bigl[p^2+(\boldsymbol{p}\cdot\boldsymbol{n})^2\bigr].
\end{equation*}

\end{xiti}
